\chapter{Cricothyrotomy}

\section*{Overview}
Cricothyrotomy is an emergency surgical procedure that opens the cricothyroid membrane to establish a definitive airway when endotracheal intubation is impossible or contraindicated. It provides direct access to the trachea for ventilation and oxygenation.

\section*{Alternative Names}
Cricothyroidotomy, needle cricothyrotomy, surgical airway, emergency tracheostomy (less precisely)

\section*{Principle}
The cricothyroid membrane lies between the thyroid and cricoid cartilages, beneath the skin of the anterior neck, and is relatively avascular and superficial. Incising or puncturing this membrane allows a tube or catheter to be placed into the subglottic trachea, bypassing the larynx and upper airway obstruction.

\section*{History}
Emergency airway access through the cricothyroid membrane has been performed for over a century and became standardized in trauma and airway management algorithms. It remains the definitive rescue airway in the "cannot intubate, cannot oxygenate" scenario.

\section*{Why It Is Done}
Performed in life-threatening airway obstruction or failure when mask ventilation and endotracheal intubation have failed or are impossible, including severe facial trauma, laryngeal obstruction, anaphylaxis, foreign body aspiration, or massive upper airway hemorrhage.

\section*{Is This a Primary or Secondary Test or Procedure}
A secondary rescue procedure reserved for emergency situations when standard airway management fails, not a primary elective airway.

\section*{How It Is Performed}
The patient is positioned with the neck extended, the laryngeal prominence is identified, and the cricothyroid membrane is palpated. In the open technique, a small vertical skin incision is made, the membrane is incised horizontally, and a small endotracheal or tracheostomy tube is guided into the trachea and secured. In the needle technique, a large-bore catheter is inserted through the membrane and connected to an oxygen source for jet ventilation.

\section*{Preparation}
No preparation is possible in the emergency setting; positioning, sterilization, and immediate availability of equipment are the priorities. Consent is implied by the life-threatening nature of the situation.

\section*{Contraindications and Precautions}
Relative contraindications include age under approximately 12 years (where needle cricothyrotomy may be preferred), laryngeal trauma or fracture that obliterates landmarks, and bleeding diathesis. Precautions include avoiding injury to the cricoid or thyroid cartilages and confirming intratracheal placement before ventilation.

\section*{Risks}
Bleeding and hematoma, damage to the larynx or trachea, subglottic stenosis, infection, subcutaneous emphysema, and esophageal or vascular injury from misplacement.

\section*{Aftercare and Recovery}
Ventilation is confirmed by chest rise and capnography, and the patient is converted to a formal tracheostomy or endotracheal airway as soon as conditions allow. The patient is monitored in an intensive care setting for bleeding, infection, and airway complications.

\section*{Results and Interpretation}
Success is confirmed by visible chest rise, exhaled carbon dioxide, and oxygen saturation improvement after connection to a ventilation circuit. Failure to ventilate indicates esophageal or tissue placement requiring immediate repositioning.

\section*{Expected Results}
A patent airway with adequate oxygenation and ventilation, stable tube position, and a subsequent planned transition to a definitive airway.

\section*{Follow-up}
Formal tracheostomy or airway reassessment, chest radiography, and clinical evaluation for laryngeal or tracheal injury. Early conversion to a definitive airway and careful decannulation planning are indicated.

\section*{References}
\begin{enumerate}
  \item MedlinePlus. Cricothyrotomy. U.S. National Library of Medicine; 2025.
  \item Current Medical Diagnosis and Treatment. McGraw-Hill; 2025.
\end{enumerate}

\clearpage
